\documentclass[11pt]{article}
\usepackage[a4paper,margin=25mm]{geometry}
\usepackage{amsmath,amssymb,amsthm,mathtools,booktabs,microtype}
\usepackage[T1]{fontenc}
\usepackage{lmodern}
\usepackage{xcolor}
\usepackage[colorlinks=true,linkcolor=blue!55!black,urlcolor=blue!55!black,citecolor=blue!55!black]{hyperref}
\usepackage{fancyhdr}
\pagestyle{fancy}\fancyhf{}\fancyhead[L]{\small Universal prime fingerprints}\fancyhead[R]{\small Research note $\cdot$ 27 November 2026}\fancyfoot[C]{\thepage}
\setlength{\headheight}{14pt}
setlength{\emergencystretch}{1em}
\newtheorem{theorem}{Theorem}[section]
\newtheorem{lemma}[theorem]{Lemma}
\newtheorem{corollary}[theorem]{Corollary}
\newtheorem{proposition}[theorem]{Proposition}
\theoremstyle{remark}\newtheorem{remark}[theorem]{Remark}
\newcommand{\F}{\mathbb F}
\newcommand{\Pp}{\mathbb P}
\newcommand{\D}{\mathcal D}
\newcommand{\C}{\mathcal C}
\newcommand{\abs}[1]{\left|#1\right|}
\newcommand{\GRH}{\textup{DGRH}}
\newcommand{\wt}{\operatorname{wt}}
\title{\bf Universal prime fingerprints at list-decoding capacity\\[4pt]\large An arithmetic coding theorem from Dirichlet GRH}
\author{Research note for the ProveIt project}
\date{27 November 2026}
\begin{document}
\maketitle
\begin{abstract}
Assuming the Dirichlet generalized Riemann hypothesis (DGRH), we prove that a single list of distinct primes, all $O((\log Q)^2)$, identifies every primitive quadratic character of conductor at most $Q$ up to a constant-size list, even in the presence of adversarial errors, erasures, and ramification. The length of the prime list approaches the exact information-theoretic capacity of this problem. The result extends to characters of order dividing any fixed prime $\ell$, with the corresponding $\ell$-ary entropy. The key step is a uniform multiplicative comparison between the law of evaluations of any bounded number of independent characters at a sampled prime and the uniform distribution. A rank-witness argument converts this local estimate into a simultaneous coding guarantee for all characters and all received words. The proof is classical and self-contained apart from standard analytic estimates, with the precise GRH dependency recorded. An exact finite example and its reproducible verifier accompany the note.
\end{abstract}

\section{Statement and scope}
Let $\ell$ be a fixed prime and let $\C_\ell(Q)$ consist of the primitive Dirichlet characters $\chi$ of conductor $f_\chi\le Q$ whose order divides $\ell$, including the character of conductor $1$. Here a character of order dividing $\ell$ takes values in the group of $\ell$th roots of unity on integers coprime to its conductor, and is zero elsewhere. Write $\zeta_\ell=\exp(2\pi i/\ell)$ and $M=\abs{\C_\ell(Q)}$.

A prime list $P=(p_1,\ldots,p_n)$ gives a fingerprint of $\chi$: at a prime not dividing $f_\chi$, its symbol is the unique $a\in\F_\ell$ with $\chi(p_i)=\zeta_\ell^a$. If $p_i\mid f_\chi$, its symbol is a wildcard, denoted by $*$. A received word consists of a set $I\subseteq\{1,\ldots,n\}$ of observed positions and symbols $y\in\F_\ell^I$. Fix an erasure fraction $e\in[0,1)$ and an error fraction $b\in[0,1-1/\ell)$. A candidate $\chi$ is compatible with $(I,y)$ if
\begin{equation}\label{eq:compatible}
\abs{I}\ge(1-e)n,\qquad
\#\{i\in I:p_i\nmid f_\chi,\ \chi(p_i)\ne\zeta_\ell^{y_i}\}\le b\abs{I}.
\end{equation}
Thus ramified positions are free mismatches, in addition to ordinary errors and erasures. The error fraction is measured relative to the number of observed positions, not the original length.

Define the $\ell$-ary entropy and the capacity by
\begin{align*}
H_\ell(u)&=-u\log_\ell u-(1-u)\log_\ell(1-u)+u\log_\ell(\ell-1),\\
K_\ell(e,b)&=(1-e)(1-H_\ell(b)),
\end{align*}
with $0\log 0=0$. All other logarithms are natural unless a base is displayed.

\begin{theorem}[Universal prime fingerprints; uses DGRH]\label{thm:main}
Fix $\ell,e,b$ as above and any rate $0<R<K_\ell(e,b)$. There are constants $B,c>0$ and an integer $r\ge1$, depending only on these fixed parameters, such that for all sufficiently large $Q$ there exist
\[
 n=\left\lceil\frac{\log_\ell M}{R}\right\rceil
\quad\text{distinct primes}\quad
 B(\log Q)^2<p_i<2B(\log Q)^2
\]
with the following simultaneous guarantee: every received word satisfying the first condition of \eqref{eq:compatible} has at most $\ell^{r-1}$ compatible candidates in $\C_\ell(Q)$. One may take any integer $r$ satisfying
\begin{equation}\label{eq:rankchoice}
 r\bigl(K_\ell(e,b)-R\bigr)>\log_\ell(\ell+1).
\end{equation}
For the weighted sampling without replacement described below, the failure probability is at most $\ell^{-cn}$.
\end{theorem}

\begin{theorem}[Converse; unconditional]\label{thm:converse}
Suppose a prime list of length $n=n(Q)$, with no restriction on the sizes of its primes, has at most $L$ compatible candidates for every received word, for a fixed $L$ independent of $Q$. Then
\[
 \limsup_{Q\to\infty}\frac{\log_\ell M}{n(Q)}\le K_\ell(e,b).
\]
The same bound holds for any alphabet-$\ell$ fingerprint scheme, even without ramification.
\end{theorem}

In the quadratic case, $\C_2(Q)$ is naturally indexed by fundamental discriminants $D$ with $0<|D|\le Q$, including $D=1$; the symbols at odd primes are Legendre symbols $(D/p)$. In particular, for any fixed $\epsilon>0$ one may take $n\le (1+\epsilon)\log_2 Q/K_2(e,b)+O(1)$ for large $Q$ by using a rate slightly greater than $K_2(e,b)/(1+\epsilon)$. Theorem~\ref{thm:converse} says that this leading constant cannot be improved if the list bound remains constant. The parameters must be fixed before $Q$ tends to infinity; no claim of uniformity as $R$ approaches capacity is made.

\section{Relation to the supplied Lean proof and article}
The vendored \texttt{grh-lean/grh} directory is documented as a byte-identical copy of the readable edition of upstream commit \texttt{72fead102f663d5c7a40}\allowbreak\texttt{7964a13dc0cf9d977207}, with Lean 4.35.0 and Mathlib commit \texttt{56a005855ea17e0b0958}\allowbreak\texttt{f73ec801ae140b0a1679} \cite{vendor}. Its \texttt{Statements.lean} defines four propositions using Mathlib, and \texttt{Main.lean} exports five theorems. The exact analytic inputs needed here are \texttt{LFunction\_ne\_zero\_}\allowbreak\texttt{of\_half\_lt\_re} and \texttt{LFunction\_eq\_zero\_re}, and their zeta counterparts (or the primitive specialization for the latter). They cover every positive modulus, with the usual pole exclusion. In particular, \texttt{New.Main} obtains the strip assertion from the one-sided nonvanishing assertion using the functional equation. The latter is obtained in \texttt{New.Reduction.}\allowbreak\texttt{OneSidedBound}, via a one-sided smoothed prime remainder estimate and a zero-detection criterion. The vendored package is separate from ProveIt's root build.

The merged article \cite{article} (723 pages according to its README) treats the source theorem as an input, rather than independently rechecking its proof. Its Section~\texttt{dir.tex} contains a uniform smoothed prime-sum lemma (label \texttt{dir:lem:smooth}), a small-bias prime pool lemma (\texttt{dir:lem:bias}), squarefree product support and relative counts (\texttt{dir:thm:quotient}, \texttt{dir:thm:short}), and a theorem on realizing most of a prescribed set of quadratic signs with one prime (\texttt{dir:thm:signs}). The present theorem reverses the latter perspective: instead of finding a prime to match a given sign word, we select one short prime list that separates \emph{all} characters up to controlled ambiguity under hostile observations. The result is not a restatement of the short-product theorems: the observations are evaluations at individual primes, not residues of their products. Sections on recovery from zeros or prime sums (\texttt{rec.tex}) and on the analytic mechanism (\texttt{src.tex}, \texttt{ana.tex}) provide related context.

Independent checks of the pinned upstream commit are documented by Sutherland \cite{sutherland}; the source README also reports standard-axiom-only dependencies and independent proof checkers. We inspected the exported statements, selected source modules and relevant article sections, not all of the approximately 300,000 lines of Lean or all 914 article results. We did not rebuild the upstream project or formalize the new theorem. We label the new implication as using DGRH, so its dependence is explicit regardless of the status of independent verification. No GRH for general Artin, Hecke or nonabelian Dedekind functions is invoked.

\section{Analytic input: nearly uniform bounded-rank evaluations}
Fix a nonnegative $C^\infty$ function $w$ of compact support in $(1,2)$ with $\int w=1$. For example, normalize $\exp(-1/((u-5/4)(7/4-u)))$ on $(5/4,7/4)$ and extend by zero. Put
\[
 W_X=\sum_p w(p/X)\log p,\qquad
 \mu_X(p)=\frac{w(p/X)\log p}{W_X}.
\]
Primes of zero weight are omitted. In the following estimates all implicit constants may depend on $w$, but not on the modulus or $X$.

\begin{lemma}[Smoothed prime sum; uses DGRH]\label{lem:prime}
For any Dirichlet character $\eta$ modulo $q$ and $X\ge2$,
\[
 \sum_p w(p/X)\log p\,\eta(p)
 =\mathbf 1_{\eta=\eta_0}X+O_w\bigl(\sqrt X\log(2q)+\sqrt X\log^2(2X)\bigr).
\]
In particular $W_X\sim X$ and $\max_p\mu_X(p)=O_w(\log X/X)$.
\end{lemma}
\begin{proof}
This is the complex-character version of the article's smoothed prime-sum lemma, followed by removal of higher prime powers. For clarity, the relevant explicit-formula argument is as follows. Let $\eta^*$ be the primitive character inducing $\eta$, of conductor $f\mid q$; take the zeta function for $f=1$. The Mellin transform $\widehat w(s)=\int_0^\infty w(u)u^{s-1}\,du$ decays faster than any power of $1+|\Im s|$ on bounded vertical strips. Mellin inversion gives
\[
 \sum_{m\ge1}\Lambda(m)\eta^*(m)w(m/X)
 =\frac1{2\pi i}\int_{(2)}-\frac{L'}{L}(s,\eta^*)X^s\widehat w(s)\,ds.
\]
Shift to $\Re s=-1/2$, avoiding zeros at the horizontal edges. The pole at $s=1$ for $\zeta$ supplies $X$; nontrivial zeros $\rho$ contribute $-X^\rho\widehat w(\rho)$ with multiplicity; the possible trivial zero at $s=0$ contributes $O_w(1)$. The standard zero count in unit height intervals is $O(\log(2f(|t|+2)))$ (see \cite{MV}); DGRH gives $\Re\rho=1/2$, and hence the zero sum is absolutely convergent and $O_w(\sqrt X\log(2f))$. The same zero count and the local logarithmic derivative formula justify taking horizontal heights away from zeros; the functional equation and the absolutely convergent Euler product on $\Re(1-s)=3/2$ bound $L'/L$ on the new line by $O(\log(2f(|t|+2)))$. Its integral is $O_w(X^{-1/2}\log(2f))$. These estimates are valid for nonreal characters as well; they do not require a symmetry of the ordinates of one character.

Inflating the primitive character to modulus $q$ changes the sum by $O_w(\log q)$: at most one power of any fixed prime lies in $(X,2X)$. Removing the terms with $m=p^k$, $k\ge2$, costs at most $O_w(\sqrt X\log^2(2X))$, by counting integers up to $(2X)^{1/k}$ and summing over $k\le\log_2(2X)$. This proves the estimate. Taking modulus one gives the claims about $W_X$ and the maximal atom.
\end{proof}

To describe independence, regard primitive characters of order dividing $\ell$ as elements of the direct limit of the character groups of $(\mathbb Z/m\mathbb Z)^\times$, with inflation along divisibility. This is an $\F_\ell$-vector space (multiplication of characters is addition, and taking powers is scalar multiplication). A primitive character has a unique conductor and is the unique primitive representative of its element in this limit. Thus $r$ such characters are independent precisely when every nontrivial product $\prod_{j=1}^r\chi_j^{a_j}$, $a_j\in\{0,\dots,\ell-1\}$, is nonprincipal after inflation to a common modulus.

For a character $\chi$, temporarily pad its wildcard by the symbol zero; call the padded evaluation $A_\chi(p)\in\F_\ell$. This padding is only a proof device; the decoder will allow all values at ramified primes.

\begin{lemma}[Bounded-rank equidistribution; uses DGRH]\label{lem:rank}
Fix $r,\ell$ and $\eta>0$. There is $B=B(w,r,\ell,\eta)>0$ such that for all sufficiently large $Q$, for $X=B(\log Q)^2$ and every independent $r$-tuple $\chi_1,\dots,\chi_r\in\C_\ell(Q)$, the law under $p\sim\mu_X$ of $(A_{\chi_1}(p),\ldots,A_{\chi_r}(p))$ satisfies, for every $v\in\F_\ell^r$,
\begin{equation}\label{eq:pointwise}
 \Pp(A=v)\le (1+\eta)\ell^{-r}.
\end{equation}
The estimate is uniform in the tuple. It holds with lower bound $(1-\eta)\ell^{-r}$ as well.
\end{lemma}
\begin{proof}
Set $m=\operatorname{lcm}(f_{\chi_1},\ldots,f_{\chi_r})\le Q^r$. For a nonzero Fourier index $a\in\F_\ell^r$, the expectation of $\zeta_\ell^{a\cdot A}$ agrees at every prime not dividing $m$ with the normalized prime sum of the nonprincipal character $\prod_j\chi_j^{a_j}$ inflated to $m$. The discrepant contribution is bounded by $\|w\|_\infty\log m/W_X$, since the inflated character vanishes at primes dividing $m$. Lemma~\ref{lem:prime} therefore bounds the absolute value of every nontrivial Fourier coefficient by
\[
 O_w\left(\frac{r\log(2Q)}{\sqrt X}+\frac{\log^2(2X)}{\sqrt X}+\frac{r\log Q}{X}\right)
 =O_w(r/\sqrt B)+o(1).
\]
Fourier inversion on $\F_\ell^r$ bounds the pointwise relative error from uniform by $(\ell^r-1)$ times this quantity. First choose $B$ sufficiently large, then $Q$ large. This yields both bounds. In particular no statistical independence of primes or of zeros has been assumed.
\end{proof}

\section{Proof of the coding theorem}
\subsection{Size of the family}
\begin{lemma}\label{lem:count}
For fixed prime $\ell$, $M=\abs{\C_\ell(Q)}=Q^{1+o(1)}$. For $\ell=2$, $M\asymp Q$.
\end{lemma}
\begin{proof}
For an odd prime power $p^a$, the group of units is cyclic, so there are at most $\ell$ characters of order dividing $\ell$. At a power of two the unit group is generated by at most two elements, giving at most $\ell^2$. By CRT the number at modulus $f$ is at most $\ell^{\omega(f)+1}$, which is $f^{o(1)}$ uniformly as $f\to\infty$: a number with $k$ distinct prime factors is at least $(k+1)!$. Summing over $f\le Q$ proves $M\le Q^{1+o(1)}$. For each prime $p\le Q$ with $p\equiv1\pmod\ell$, the cyclic unit group has $\ell-1$ nontrivial order-$\ell$ characters and all are primitive. The prime number theorem in this fixed progression gives $M\gg_\ell Q/\log Q$ (for $\ell=2$ use all odd primes). This yields the first assertion unconditionally.

For $\ell=2$, the description by fundamental discriminants gives $M\le2Q+1$. A positive proportion of integers $D\le Q$ with $D\equiv1\pmod4$ are squarefree: the number removed by odd prime squares is at most $(Q/4)\sum_{p\ge3}p^{-2}+O(\sqrt Q)$, and the sum is less than $1/2$. All the remaining $D$ are fundamental. This proves $M\asymp Q$.
\end{proof}

\subsection{Sampling without replacement}
Choose the $n$ primes successively from $\mu_X$, at each step renormalizing its restriction to primes not yet drawn. Lemma~\ref{lem:count} gives $n=O(\log Q)$. By Lemma~\ref{lem:prime}, the total mass removed at any step is at most
\[
 (n-1)\max_p\mu_X(p)=O\left(\frac{\log\log Q}{\log Q}\right)=o(1).
\]
In particular the selection can always proceed. Given any fixed independent $r$-tuple and any past draws, the probability that the next column of its padded evaluations equals a specified $v$ is at most $(1+\eta_0)\ell^{-r}/(1-o(1))$, where $\eta_0$ is the tolerance in Lemma~\ref{lem:rank}. For any prescribed $\eta>0$, choose $\eta_0<\eta$ and enlarge $Q$ so that this bound is at most $(1+\eta)\ell^{-r}$. Consequently, for any fixed subset $I$ of $m$ positions and any prescribed $r\times m$ matrix of symbols there, the probability of this matrix is at most
\begin{equation}\label{eq:matrix}
 (1+\eta)^m\ell^{-rm}.
\end{equation}
This follows by conditioning in chronological order on all earlier draws, including those outside $I$; no independence of columns after conditioning is required.

\subsection{Rank witnesses and Hamming balls}
Write $K=K_\ell(e,b)$, and choose $r$ satisfying \eqref{eq:rankchoice}. Let
\[
 d=r(K-R)-\log_\ell(\ell+1)>0.
\]
By continuity of entropy, choose $b'$ with $b<b'<1-1/\ell$ so close to $b$ that
\begin{equation}\label{eq:slack}
 r\bigl(K-(1-e)(1-H_\ell(b'))\bigr)<d/3.
\end{equation}
Choose $\eta>0$ with $\log_\ell(1+\eta)<d/3$, and choose $B$ as in the preceding sampling argument.

Every character of conductor $f\le Q$ has at most $\log Q/\log X$ ramified positions among the selected distinct primes: their product divides $f$. For $m=|I|\ge(1-e)n$, this number is $o(m)$, uniformly in $f$, because $n\asymp\log Q$ and $\log X\to\infty$. For sufficiently large $Q$ it is less than $(b'-b)m$. Hence every compatible candidate has a padded fingerprint in the $\ell$-ary Hamming ball of radius $\lfloor b'm\rfloor$ about $y$ on $I$.

If more than $\ell^{r-1}$ candidates exist, their span in the direct-limit vector space has dimension at least $r$: a vector space of dimension at most $r-1$ over $\F_\ell$ has at most $\ell^{r-1}$ elements. Thus some independent $r$ of the candidates witness failure. This argument includes the principal character (the zero vector) and does not assume all candidates independent.

The volume of a Hamming ball is
\[
 V_\ell(m,k)=\sum_{j=0}^k\binom mj(\ell-1)^j,
\qquad
 V_\ell(m,\lfloor b'm\rfloor)\le\ell^{mH_\ell(b')}.
\]
For the inequality, put $z=b'/((\ell-1)(1-b'))\in(0,1)$ and use $z^j\ge z^{b'm}$ for $j\le b'm$ in the binomial expansion of $(1+(\ell-1)z)^m$. This also proves the limiting case $b'=0$ directly if needed.

For fixed $I,y$ and fixed independent $r$-tuple, there are at most $V_\ell(m,\lfloor b'm\rfloor)^r$ possible evaluation matrices placing all $r$ padded words in the ball. By \eqref{eq:matrix} their combined probability is at most
\[
 (1+\eta)^m\ell^{-rm(1-H_\ell(b'))}.
\]
There are at most $M^r$ ordered tuples and at most $(\ell+1)^n$ received words over all erased sets (each position is erased or one of $\ell$ symbols). Using $m\ge(1-e)n$ and $\log_\ell M\le Rn$, the union bound gives
\begin{align*}
 \Pp(\text{failure})
 &\le M^r(\ell+1)^n(1+\eta)^n
   \ell^{-rn(1-e)(1-H_\ell(b'))}\\
 &\le \ell^{n[\log_\ell(\ell+1)+rR+\log_\ell(1+\eta)-r(1-e)(1-H_\ell(b'))]}\\
 &\le \ell^{-dn/3}.
\end{align*}
The last inequality uses \eqref{eq:slack} and the choice of $\eta$. This is strictly less than one. It proves Theorem~\ref{thm:main}, including the high-probability assertion with $c=d/3$. \qed

\section{Information-theoretic converse and consequences}
\begin{proof}[Proof of Theorem~\ref{thm:converse}]
Fix a single observed set $I$ of size $m=\lceil(1-e)n\rceil$ and pad every wildcard with zero. A Hamming ball of radius $k=\lfloor bm\rfloor$ about any padded fingerprint is contained in the set of received words compatible with that character. Double-count pairs $(\chi,y)$ where $y$ lies in this ball. Counting by characters gives $MV_\ell(m,k)$; counting by received words gives at most $L\ell^m$. Thus
\begin{equation}\label{eq:converse}
 \log_\ell M\le m-\log_\ell V_\ell(m,k)+\log_\ell L.
\end{equation}
Even the weaker inequality $M\le L\ell^n$ forces $n\to\infty$. For $b>0$, Stirling's formula applied to the $k$th term and the preceding upper bound give $\log_\ell V_\ell(m,k)=mH_\ell(b)+O(\log m)$; for $b=0$ the identity is exact. Divide \eqref{eq:converse} by $n$, use $m/n\to1-e$, and take the limsup. No facts about primes enter this bound.
\end{proof}

\begin{corollary}[Quadratic characters at the optimal leading length]\label{cor:quad}
For every $e<1$, $b<1/2$, and every $\epsilon>0$, there is a constant $L$ and, for each sufficiently large $Q$, a list of at most
\[
 (1+\epsilon)\frac{\log_2 Q}{(1-e)(1-H_2(b))}+O(1)
\]
distinct primes of size $O((\log Q)^2)$, for which every allowed observation of Legendre symbols $(D/p)$ is compatible with at most $L$ fundamental discriminants $|D|\le Q$. The leading coefficient is optimal in the constant-list regime. A single list works for both positive and negative discriminants and does not require the conductor to be known.
\end{corollary}
\begin{proof}
Apply the main theorem and Lemma~\ref{lem:count} with $\ell=2$, choosing $R$ between $K_2(e,b)/(1+\epsilon)$ and $K_2(e,b)$. The converse yields optimality.
\end{proof}

For example, for binary symbols, $e=0.1$ and $b=0.1$ give $K_2\simeq0.4779$. The optimal leading length is approximately $2.092\log_2 Q$; at rate $R=0.4$ the theorem gives length $2.5\log_2 Q+O(1)$ and one may take $r=21$ (thus a list bound $2^{20}$). This bound is intentionally elementary rather than practically optimized; stronger finite-rank random-code arguments may reduce the constant list size. The prime height $B$ in our proof is also not optimized: it depends on the tolerated rank-$r$ Fourier error, and may be very large. A trivial necessary height scale is $\gg\log Q\log\log Q$ for $\asymp\log Q$ distinct prime observations, by the prime number theorem. Closing the gap to the sufficient $O((\log Q)^2)$ is a different problem.

\section{Finite check and reproducibility}
The accompanying standard-library Python script exhaustively checks a small instance, without relying on DGRH. Set $Q=128$, take 16 distinct odd primes between 101 and 300 as listed below, and allow at most two erased positions and two errors among the surviving positions (absolute error count for this experiment). Wildcards at ramified positions are permitted. The code enumerates all fundamental discriminants, each allowed erased set, and every word in the wildcard-expanded ball around each candidate. The supports of wildcard flips and error flips are disjoint, so a character contributes at most once per received word. The exact maximum list size is five.
\begin{center}
\begin{tabular}{rrrrrrrr}
\toprule
103&113&139&157&163&167&179&181\\
197&199&223&233&251&271&277&293\\
\bottomrule
\end{tabular}
\end{center}
\begin{center}
\begin{tabular}{rr}
\toprule
List size & Number of pairs (erased set, received word)\\
\midrule
0 & 1,588,406\\
1 & 772,824\\
2 & 170,672\\
3 & 22,278\\
4 & 1,655\\
5 & 69\\
\bottomrule
\end{tabular}
\end{center}
There are 79 characters in this example. One maximum witness erases the first two primes and uses the signs
\[
 (-,-,+,+,-,-,+,+,-,-,-,-,-,+),
\]
with candidate discriminants $-83,24,61,97,104$. The script records the full histogram and this witness in \texttt{example.json}. This is an exact finite verification, not a test of the asymptotic constants or the analytic input.

\section{Novelty and limits}
The probabilistic mechanism is related to random linear code list decoding: in an ordinary random linear code, independent messages give independent uniformly distributed codewords. Guruswami, H\aa stad and Kopparty \cite{GHK} prove much sharper list-size bounds for genuinely random linear codes. The present model is different: the messages are primitive characters of bounded conductor, not a fixed full vector space, and each coordinate must be an arithmetic evaluation at a distinct small prime. The uniform bounded-rank comparison, the wildcard allowance for all conductors, and the capacity-achieving arithmetic restriction are the proposed contribution. The elementary rank argument yields a constant list bound exponential in the inverse rate gap, not the best random-code list size.

The supplied article sections reviewed and a targeted literature search did not reveal this exact formulation. This is a candidate novel theorem, not a certified priority claim. In particular, the proof deliberately uses the established random-code rank-witness method rather than claiming it as new. It gives existence and high-probability randomized selection, not an efficient worst-case decoder or an explicit numerical constant $B$. Exhaustive decoding by enumerating conductors and characters is always possible; the finite experiment implements it for quadratic characters. No claim is made that all 131 source reports have been searched exhaustively, that the analytic bridge is already in Lean, or that this note is independently refereed.

{\small\begin{thebibliography}{9}\setlength{\itemsep}{2pt}
\bibitem{vendor} ProveIt, \emph{Vendored grh-lean (readable edition)}, November 2026. \href{https://github.com/VladimirReshetnikov/ProveIt/tree/main/Analysis/GeneralizedRiemannHypothesis/vendor/grh-lean}{Repository directory}; upstream \href{https://github.com/lsernau/grh-lean/tree/72fead102f663d5c7a407964a13dc0cf9d977207/grh}{pinned edition}.
\bibitem{article} ProveIt Contributors, \emph{Beyond the Riemann hypothesis for Dirichlet $L$-functions: consequences, transfers and limits}, 26 November 2026. \href{https://github.com/VladimirReshetnikov/ProveIt/tree/main/Analysis/GeneralizedRiemannHypothesis/Article}{Article sources}, especially \texttt{sections/dir.tex}.
\bibitem{sutherland} A. V. Sutherland, \emph{Verification files for the shorter Lean proof of RH and GRH}, 26 November 2026. \url{https://klein.mit.edu/~drew/shortgrh/}.
\bibitem{MV} H. L. Montgomery and R. C. Vaughan, \emph{Multiplicative Number Theory I: Classical Theory}, Cambridge University Press, 2007; Dirichlet $L$-function zero counting, explicit formulae and conditional prime sums, Chapter 13 and preceding background.
\bibitem{GHK} V. Guruswami, J. H\aa stad and S. Kopparty, On the list-decodability of random linear codes, \emph{IEEE Trans. Inform. Theory} 57 (2011), 718--725. \url{https://arxiv.org/abs/1001.1386}.
\end{thebibliography}}
\end{document}
